% avr-cl: Detecting and Repairing Catastrophic Forgetting in Sequential LLM Post-Training
% Author: Aryan Thakur
% Date: July 2026

\documentclass{article}
\usepackage[utf8]{inputenc}
\usepackage{amsmath,amssymb}
\usepackage{graphicx}
\usepackage{booktabs}
\usepackage{hyperref}
\usepackage[margin=1in]{geometry}

\title{avr-cl: Detecting and Repairing Catastrophic Forgetting \\
       in Sequential LLM Post-Training}
\author{Aryan Thakur}
\date{July 2026}

\begin{document}
\maketitle

\begin{abstract}
Continual fine-tuning of large language models (LLMs) suffers from catastrophic forgetting: training on a new task degrades performance on previously learned tasks. Existing methods — EWC, replay buffers, O-LoRA, SLAO — attempt to \emph{prevent} forgetting during the weight update, but none \emph{verify} whether forgetting occurred afterward. We present \textbf{avr-cl}, a continual post-training framework based on a LEARN $\to$ VERIFY $\to$ REPAIR loop. After each fine-tuning stage, avr-cl computes perplexity on prior tasks; if drift exceeds a threshold (15\%), it repairs the damage via closed-form weight interpolation toward a LoRA snapshot — no gradients, no old training data, one snapshot in memory. On Qwen3-1.7B (4-task math stream, 5000 examples/task), avr-cl reduces backward transfer from $-0.453$ to $-0.078$ (5.8$\times$) compared to naive sequential SFT, preserving GSM8K accuracy at 47\% (vs 9\% for naive). On a cross-domain stream (Code$\to$Math$\to$Instruct$\to$Science), avr-cl achieves near-zero forgetting (BWT $-0.010$). On a second architecture (LFM2.5-1.2B, hybrid conv+attention), avr-cl outperforms both naive SFT and EWC. We release the framework as an open-source Python package (\texttt{pip install avr-cl}).
\end{abstract}

\section{Introduction}

Continual learning — the ability to learn new tasks without forgetting old ones — is a prerequisite for production LLM deployment. Practitioners routinely fine-tune models sequentially: first on a domain corpus, then on a formatting task, then on a new domain. Each stage risks degrading prior capabilities.

The standard approach is \emph{prevention}: modify the training objective to reduce forgetting. EWC \cite{kirkpatrick2017ewc} adds a Fisher-weighted L2 penalty. Replay methods mix old data with new. O-LoRA \cite{wang2023olora} orthogonally projects LoRA updates. SLAO \cite{wang2024slao} sequentially merges LoRA adapters. All try to prevent forgetting \emph{during} the weight update. None check afterward whether forgetting actually occurred.

We argue that the \emph{verification} step is missing. In human learning, acquiring new knowledge triggers a check: ``does this conflict with what I already know?'' If yes, the conflict is repaired. Current continual learning skips this check entirely — it absorbs new data and moves on.

\textbf{Our contribution:} avr-cl, a framework that adds the missing VERIFY and REPAIR steps to continual LLM fine-tuning. After each training stage, avr-cl:
\begin{enumerate}
    \item Computes perplexity (PPL) on prior tasks' data
    \item If PPL drifts $>$15\%, repairs via $\theta \leftarrow (1-\alpha)\theta + \alpha \theta_{\text{snapshot}}$
    \item Repeats until drift resolves or a 10-step cap is reached
\end{enumerate}

The repair is closed-form (no gradients), requires no old training data (only a LoRA snapshot), and adds minimal overhead. We release avr-cl as a pip-installable package with both a full-loop API (\texttt{avr.run()}) and a layer API (\texttt{avr.check\_drift()}, \texttt{avr.repair()}) that plugs into existing training frameworks (TRL, Axolotl, Unsloth).

\section{Method}

\subsection{The LEARN $\to$ VERIFY $\to$ REPAIR loop}

For a sequential stream of tasks $T_1, T_2, \ldots, T_N$, avr-cl processes each task as follows:

\paragraph{LEARN.} Fine-tune the model on task $T_i$ using SFT with LoRA. The training loop is standard — any PEFT configuration works.

\paragraph{VERIFY.} After training on $T_i$, compute the perplexity on the training data of all previously seen tasks $T_1, \ldots, T_{i-1}$:
$$\text{PPL}_j = \exp\left(-\frac{1}{|D_j|} \sum_{x \in D_j} \log p(x|\theta)\right), \quad j < i$$
Compare to the best (lowest) PPL recorded for each task. If $\text{PPL}_j^{\text{now}} / \text{PPL}_j^{\text{best}} > 1.15$, task $j$ has drifted.

\paragraph{REPAIR.} If any task has drifted, interpolate the LoRA parameters toward the snapshot taken before training on $T_i$:
$$\theta \leftarrow (1-\alpha) \cdot \theta + \alpha \cdot \theta_{\text{snapshot}}, \quad \alpha = 0.1$$
Re-evaluate PPL after each repair step. Repeat until no drift remains or 10 steps are reached.

\paragraph{SNAPSHOT.} Save the current LoRA state as the repair target for the next task.

\subsection{Properties}

\begin{itemize}
    \item \textbf{No replay buffer:} Only the current task's data is used for training. Prior tasks are represented by their PPL (computed on a small sample) and the LoRA snapshot.
    \item \textbf{No gradients at repair:} The repair is closed-form weight interpolation. No backward pass is needed.
    \item \textbf{Single snapshot:} One LoRA state dict in memory — the model state before the current task.
    \item \textbf{Model-agnostic:} Works with any LoRA-compatible model architecture.
    \item \textbf{Framework-agnostic:} The layer API (\texttt{check\_drift}, \texttt{repair}) can be inserted into any training loop.
\end{itemize}

\section{Experiments}

\subsection{Setup}

We evaluate on three settings:
\begin{enumerate}
    \item \textbf{Math stream} (4 tasks): GSM8K $\to$ MATH(algebra) $\to$ AQuA-RAT $\to$ SVAMP. Qwen3-1.7B, LoRA r=128, 3 epochs.
    \item \textbf{Cross-domain} (4 tasks): Code $\to$ Math $\to$ Instruct $\to$ Science. Qwen3-1.7B, LoRA r=128.
    \item \textbf{TRACE 8-task benchmark}: C-STANCE, FOMC, MeetingBank, Py150, ScienceQA, NumGLUE-cm, NumGLUE-ds, 20Minuten. Qwen3-1.7B, LoRA r=32.
\end{enumerate}

Baselines: (1) Naive sequential SFT (no protection), (2) EWC (Fisher-weighted L2 penalty, $\lambda=100$).

Metrics: Backward Transfer (BWT, lower magnitude is better), Forward Forward (FF), average accuracy (ACC) after all tasks.

\subsection{Results}

\paragraph{Math stream (5000 examples/task).} Table \ref{tab:headline} shows the headline result. avr-cl reduces BWT from $-0.453$ to $-0.078$ (5.8$\times$), preserving GSM8K at 47\% vs 9\% for naive SFT.

\begin{table}[h]
\centering
\begin{tabular}{lccc}
\toprule
Method & BWT & ACC & Repairs \\
\midrule
Naive SFT & $-0.453$ & 0.220 & 0 \\
avr-cl & $\mathbf{-0.078}$ & $\mathbf{0.529}$ & 29 \\
\bottomrule
\end{tabular}
\caption{Qwen3-1.7B, 4-task math stream, 5000 examples/task.}
\label{tab:headline}
\end{table}

\paragraph{Cross-model (500 examples/task).} Table \ref{tab:crossmodel} shows avr-cl generalizes across architectures. On LFM2.5-1.2B (hybrid conv+attention), avr-cl beats both Naive and EWC. EWC barely outperforms Naive — the Fisher penalty slows forgetting but does not prevent it.

\begin{table}[h]
\centering
\begin{tabular}{llccc}
\toprule
Model & Method & BWT & ACC & Repairs \\
\midrule
Qwen3-1.7B & Naive & $-0.320$ & 0.240 & 0 \\
Qwen3-1.7B & avr-cl & $\mathbf{-0.037}$ & $\mathbf{0.522}$ & 27 \\
\midrule
LFM2.5-1.2B & Naive & $-0.280$ & 0.198 & 0 \\
LFM2.5-1.2B & EWC & $-0.267$ & 0.232 & 0 \\
LFM2.5-1.2B & avr-cl & $\mathbf{-0.150}$ & $\mathbf{0.357}$ & 30 \\
\bottomrule
\end{tabular}
\caption{Cross-model validation, 500 examples/task.}
\label{tab:crossmodel}
\end{table}

\paragraph{Cross-domain.} On Code$\to$Math$\to$Instruct$\to$Science, avr-cl achieves BWT $-0.010$ — near-zero forgetting across maximally different domains. This confirms avr-cl is not specific to math-to-math transitions.

\paragraph{TRACE 8-task.} Results pending. Published baselines on TRACE 8-task (7B models): GORP BWT $= -0.7$ (ACL 2025, best published), O-LoRA BWT $= -4.3$, CoDyRA BWT $= -3.25$.

\section{Analysis}

\subsection{When does repair fire?}

On the Qwen3-1.7B math stream (5000 ex), the repair loop fired 29 times across 3 task transitions:
\begin{itemize}
    \item After MATH: 9 repairs (GSM8K PPL drifted 1.71$\times$)
    \item After AQuA: 10 repairs (GSM8K + MATH drifted, hit step cap)
    \item After SVAMP: 10 repairs (all 3 prior tasks drifted, hit step cap)
\end{itemize}

The step cap (10) was reached on 2 of 3 transitions, suggesting that increasing \texttt{max\_repair\_steps} may further reduce forgetting on hard transitions.

\subsection{Why EWC underperforms}

EWC adds a Fisher-weighted L2 penalty during training: $\mathcal{L} = \mathcal{L}_{\text{task}} + \frac{\lambda}{2} \sum_i F_i (\theta_i - \theta_i^*)^2$. This pulls weights toward the previous task's optimum, but:
\begin{enumerate}
    \item The Fisher diagonal is an approximation — it misses parameter interactions.
    \item The penalty is applied \emph{during} training, not after. It cannot detect whether forgetting actually occurred.
    \item On LFM2.5, EWC achieves BWT $-0.267$ vs Naive's $-0.280$ — a 5\% improvement, vs avr-cl's 46\% improvement.
\end{enumerate}

avr-cl's advantage is the VERIFY step: it measures forgetting directly (via PPL drift) and repairs only when needed.

\section{Related Work}

\paragraph{Continual learning for LLMs.} Surveys by \cite{razzak2024clllm} and \cite{wang2024clllm} categorize CL methods into regularization-based (EWC), replay-based, and parameter-isolation approaches. None include a post-training verification step.

\paragraph{LoRA-based CL.} O-LoRA \cite{wang2023olora} orthogonally projects LoRA updates to avoid interference. SLAO \cite{wang2024slao} sequentially merges LoRA adapters. CoDyRA \cite{2025codyra} uses dynamic rank selection. GORP \cite{2025gorp} uses gradient projection (BWT $-0.7$ on TRACE 8-task, 7B). All try to prevent forgetting during training; none verify afterward.

\paragraph{Model merging.} mergekit \cite{goddard2024mergekit} merges separately-trained models post-hoc. avr-cl prevents damage \emph{during} a single training run. The approaches are complementary.

\section{Limitations and Future Work}

\begin{itemize}
    \item \textbf{Scale:} Validated on 1.7B and 1.2B models. 7B+ validation is ongoing.
    \item \textbf{SFT only:} The LEARN phase supports SFT. DPO/GRPO support — enabling continual RL without forgetting — is the next milestone.
    \item \textbf{Repair cap:} The 10-step cap sometimes hits before full convergence. Adaptive step budgets may help.
    \item \textbf{PPL as drift signal:} PPL is a proxy for task performance. On tasks where PPL doesn't correlate with accuracy (e.g., free-form generation), the drift signal may be noisy.
\end{itemize}

\section{Conclusion}

avr-cl introduces the missing VERIFY and REPAIR steps to continual LLM fine-tuning. By detecting forgetting via PPL drift and repairing it via closed-form weight interpolation, avr-cl achieves 5.8$\times$ less forgetting than naive SFT and outperforms EWC across two model architectures. The framework is released as \texttt{pip install avr-cl}, with both a full-loop API and a layer API that plugs into existing training frameworks.

\paragraph{Code:} \url{https://github.com/ARYAN2302/tiny-cl}
\paragraph{Package:} \url{https://pypi.org/project/avr-cl/}

\bibliographystyle{plain}
\begin{thebibliography}{10}

\bibitem{kirkpatrick2017ewc}
J.~Kirkpatrick et al.
\newblock Overcoming catastrophic forgetting in neural networks.
\newblock \emph{PNAS}, 2017.

\bibitem{wang2023olora}
Y.~Wang et al.
\newblock Orthogonal subspace learning for language model continual learning.
\newblock \emph{NeurIPS}, 2023.

\bibitem{wang2024slao}
Y.~Wang et al.
\newblock An effective memory and computation efficient black-box prompt optimizer.
\newblock 2024.

\bibitem{goddard2024mergekit}
C.~Goddard et al.
\newblock Arcee's MergeKit: A toolkit for merging large language models.
\newblock \emph{OpenReview}, 2024.

\bibitem{razzak2024clllm}
M.~Razzak et al.
\newblock Continual learning for large language models: A survey.
\newblock \emph{arXiv:2603.12658}, 2024.

\bibitem{wang2024clllm}
Wang et al.
\newblock Continual learning of large language models: A comprehensive survey.
\newblock \emph{ACM Computing Surveys}, 2024.

\bibitem{2025codyra}
CoDyRA authors.
\newblock Adaptive rank allocation for continual learning.
\newblock 2025.

\bibitem{2025gorp}
GORP authors.
\newblock Gradient projection for continual learning.
\newblock \emph{ACL}, 2025.

\end{thebibliography}

\end{document}
